\section{The confidentiality claim and public view}

Qin et al.'s GBFPlus protocol records cross-domain data transfers in published
ledger attestations while supporting authorized tracing~\cite{qin2026gbfplus}.
The source article lists the adjacent domain among the protected metadata
(Table~1 and Section~4.2); Section~6.2 specifically describes recovering its
identifier only after the consistency-prefix and assigned-domain checks.
We identify an allowed execution in which an ordinary ledger observer recovers
that identifier without either check: a time window closes after one successful
insertion, and the publicly occupied cells contain exactly that insertion's
shares. XORing their payloads reconstructs the inserted word, whose suffix
identifies the adjacent domain.

Three features of the published specification suffice for this disclosure.
First, Algorithm~4 encodes an unmasked adjacent-domain identifier in the shared
word. Second, Section~5.2 publishes the finalized filter and its successful
insertion count in the transaction following source Equation~(5); Sections~5.2
and~6.2 make occupied cells distinguishable from empty cells. Third,
Section~5.2 finalizes a filter when its time window ends, without specifying
a minimum insertion count. Thus an honest administrator can initialize a
directional filter, insert one transfer, receive no further transfer in that
direction, and publish the filter when the window closes. The source's
honest-administrator model (Section~4.3) and these rules admit this execution;
they do not establish how often it occurs in deployment.

We consider one honestly generated, authentically committed snapshot.
Here, ``public'' means visible to a reader of the permissioned ledger (source
Sections~4.1--4.3). The observer needs ledger access but no regulator secret
or query token. It need not identify the publishing administrator. Repeated
reads expose the same frozen array. Authenticity and a well-formed count
alone would not establish an honest insertion history; that history is an
explicit premise here.

\section{Exact disclosure and its scope}

A filter has $m$ cells. An occupied cell $j$ stores a $\lambda$-bit word $Z_j$;
its occupancy marker distinguishes even an all-zero word from an empty cell.
Let $S_i$ be the set of $\kappa$ distinct positions selected by successful
insertion $i$. Write $N_i$ for its $\ell$-bit adjacent-domain identifier and
$P_i$ for its $(\lambda-\ell)$-bit secret-derived hash prefix, where
$0<\ell<\lambda$. These widths (with $\ell=l_{\mathrm{DA}}$ in the source)
and $\kappa$ are public configuration
parameters in source Section~5.1, Equation~(3), so the observer knows them.
Algorithm~4 enforces
\begin{equation}
  \bigoplus_{j\in S_i} Z_j = T_i,
  \qquad T_i=P_i\mathbin\Vert N_i,
  \qquad |S_i|=\kappa,
  \label{eq:row}
\end{equation}
where $\Vert$ denotes concatenation. Registered domain identifiers and their
associations are public in source Section~5.1.

The filter starts empty. A successful insertion reserves one empty selected
cell, fills the other empty selected cells with independent uniform words,
and sets the reserved word to the XOR residual required by~\eqref{eq:row}.
Already occupied words are reused and never overwritten. An insertion whose
selected cells are all occupied fails without changing the array. Consequently,
after $c$ successful insertions, the public occupied set is
$O=\bigcup_{i=1}^{c}S_i$ and every successful row equation still holds.
Let $\suffix_\ell(x)$ denote the final $\ell$ bits of $x$.

\begin{proposition}[Singleton disclosure]
\label{prop:singleton}
Every honest snapshot with $c=1$ reveals its adjacent-domain identifier:
\begin{equation}
  \suffix_\ell\!\left(\bigoplus_{j\in O}Z_j\right)=N_1.
  \label{eq:singleton}
\end{equation}
This equality holds for every selected position set, prefix, and share
realization admitted by the construction.
\end{proposition}

\begin{proof}
For the first successful insertion, all selected positions were empty and
all become occupied, so $O=S_1$. Their XOR is $T_1=P_1\Vert N_1$
by~\eqref{eq:row}; taking its suffix gives $N_1$.
\end{proof}

The observer therefore needs only the occupied words and the suffix width.
It does not have to guess the data identifier, the prefix, or the selected
positions. Figure~\ref{fig:example} distinguishes an occupied zero word from
an empty cell: the zero word counts toward occupancy but does not change the XOR.
A scan takes
$O(m+\kappa\lambda)$ elementary bit operations and
$O(\lambda+\log m)$ auxiliary bits for a singleton; ledger transfer,
parsing, and authentication are outside this scan cost.

\begin{figure}[ht]
\centering
\begin{tabular}{c*{7}{c}}
\toprule
Cell & 0 & 1 & 2 & 3 & 4 & 5 & 6\\
\midrule
Payload & -- & $010$ & -- & $111$ & -- & $000$ & --\\
\bottomrule
\end{tabular}
\[
  010\xor111\xor000=101=1\mathbin\Vert01.
\]
\caption{A singleton with three 3-bit shares and a 2-bit domain identifier.
An empty cell is marked ``--''; $000$ is an occupied payload. Public occupancy
selects all three shares, so the observer reads the identifier $01$ from their
XOR. The prefix $1$ need not be known or validated.}
\label{fig:example}
\end{figure}

The disclosure is compatible with uniform individual shares. Conditional on
any honestly reached singleton position set and prefix, and for $\kappa\geq2$,
each free word is uniform, and the residual is uniform because its XOR includes
at least one independent uniform word. Nevertheless, the complete array
always satisfies~\eqref{eq:singleton}. Two different identifiers therefore
give disjoint sets of possible complete views, even though each individual
share has the same distribution.

\paragraph{Which confidentiality claim is challenged?}
Consider a publisher with two eligible neighboring domains distinct from it.
Fix the public registration table. Internally fix the same publisher, its
secret state, data identifier, direction, window, and configuration, and choose
either neighboring identifier with equal probability for a singleton insertion.
The observer is not given the publisher's identity or secret state. In source
Algorithm~4, the adjacent identifier affects the target suffix but not the
derived positions or prefix. Those quantities can therefore be the same in
both cases, while the free shares are sampled in each case.
Equation~\eqref{eq:singleton} identifies the choice with probability one.
This comparison concerns one snapshot and its authenticating transaction;
digests, signatures, and other ledger records need not be identical between cases.

The recovered identifier names the neighbor encoded in that attestation.
It does not in general identify the anonymous publisher, associate the transfer with a
particular data identifier, or reconstruct a complete named edge or path.
This is also distinct from the source's $G_{\mathrm{qpriv}}$ game and
Theorem~4, which concern distinguishing data identifiers. Separately, source
Equation~(27) analyzes guessed tuples passing prefix and domain-membership
checks under its random-share model. The observer here performs neither
validation test: it reads a suffix from a complete share set already identified
by occupancy. Those analyses do not establish confidentiality of this suffix.

\paragraph{A limited extension to larger batches.}
Let $w=|O|$ be the public occupied-cell count. The same scan has a useful
sufficient condition for recovering an aggregate.

\begin{proposition}[Disjoint-position certificate]
\label{prop:disjoint}
For an honest snapshot with $c\geq1$, if $w=c\kappa$, then
\begin{equation}
 \suffix_\ell\!\left(\bigoplus_{j\in O}Z_j\right)
       =\bigoplus_{i=1}^{c}N_i.
 \label{eq:aggregate}
\end{equation}
\end{proposition}

\begin{proof}
Equality in
$|\bigcup_i S_i|\leq\sum_i|S_i|=c\kappa$ implies that the selected position
sets are pairwise disjoint. XORing the occupied words therefore adds each
row equation once, giving $\bigoplus_i T_i$. Taking the suffix commutes
with XOR and yields~\eqref{eq:aggregate}.
\end{proof}

For $c>1$, this is an aggregate, not a reconstruction of the individual
identifiers or their multiset. The condition requires $c\kappa\leq m$;
otherwise it is impossible. Its informational value also depends on the
candidate histories. For example, in a reachable disjoint history, changing
one eligible adjacent identifier while fixing the others changes the aggregate
and is distinguishable. If all admissible histories already have the same
aggregate, the statistic adds no information. The attack abstains when
$w\ne c\kappa$; no distribution or frequency of disjoint batches is assumed.

\section{Reproducibility checks}

The proofs establish exact identities. Supplement~S1 checks a reference
implementation of the stored-array update and the public scan; it does not
implement cryptographic position derivation, pairings, signatures, ledger
operations, or finalization. Publication and visibility premises are checked
against the source specification, rather than inferred from simulation.

The small enumerations fix position sets and prefixes and exhaust the stated
suffix/free-share choices. The larger checks construct 64 separate singleton
arrays using seeds $0$--$63$ of Python's general-purpose PRNG. These seeds
make the arrays reproducible; they do not model cryptographic randomness or
an operational workload. Table~\ref{tab:checks} states the tested equality
and input class for each row. A mismatch means the returned suffix differs
from the encoded identifier or aggregate, as applicable.

\begin{table}[ht]
\centering
\small
\begin{tabular}{@{}>{\raggedright\arraybackslash}p{0.72\linewidth}rr@{}}
\toprule
Check and input class & Cases & Mismatches\\
\midrule
Singleton: returned suffix $=N_1$; $c=1$, $\kappa=3$,
$\lambda=3$, $\ell=2$. All 4 suffixes and $8^2$ free-word pairs.
& 256 & 0\\[3pt]
Disjoint pair: returned suffix $=N_1\xor N_2$; $c=2$, $\kappa=2$,
$\lambda=3$, $\ell=2$. All $4^2$ suffix pairs and $8^2$ free-word pairs.
& 1,024 & 0\\[3pt]
Seeded singleton: returned suffix $=N_1$; $c=1$, $\kappa=12$,
$\lambda=32$, $\ell=20$, $m=64$.
& 64 & 0\\
\bottomrule
\end{tabular}
\caption{Finite implementation checks, not estimates of disclosure probability.
The small cases use $m=7$; fixed positions and prefixes are specified in
Supplement~S1. The seeded case uses the source article's Section~6.3 word widths
and $\kappa$, with illustrative $m=64$.}
\label{tab:checks}
\end{table}

Tests also check overlapping insertions, preservation of earlier row equations,
atomic failure, position order, occupied zero words, and frozen repeated reads.
The supplement supplies all arrays and a standalone command,
\texttt{python3 -S reproduce.py}, verified with Python~3.12.3.

\section{Implications and existing protections}

The protocol-specific issue is that public occupancy identifies a complete
set of shares in an allowed partial batch. The underlying distinction between
individual and joint linear observations is established; Silva and
Kschischang's coset-encoding analysis provides a general precedent
\cite{silva2011universal}. Dong, Chen, and Wen already distinguish value
disclosure from false-positive acceptance: an absent query selecting a stored item's
positions can reveal the item while returning a negative membership result
(Section~4.1)~\cite{dong2013gbf}. Here, public occupancy identifies a complete
share set without a query or a collision. Their original garbled Bloom filter
also fills unused cells with random words (Algorithm~1, lines~22--26), and their
private-set-intersection protocol uses oblivious transfer to supply an
intersection structure.
That protocol does not give the same unauthorized observer the full server
filter together with genuine insertion occupancy.

For the all-cell XOR alone, one uniform padding word independent of the remaining
cells makes the statistic uniform. An observer who retains the genuine occupancy selector
can still omit the padding and recover the original occupied-cell XOR.
For a singleton or a disjoint batch, fresh independent uniform secret masks on the
adjacent identifiers make the recovered suffix uniform, provided the mask vector
is independent of all identifiers and prefixes. Secret PRF
masking of arbitrary values is already used in oblivious key-value
retrieval~\cite{hao2024okvr}. Neither observation proves
a complete repair of GBFPlus. A repair must account for the entire public view,
authorized recovery, and key and mask reuse across releases.

A minimum batch size above one excludes the singleton execution. It excludes
the disjoint-position certificate only if accepted batches cannot satisfy
$w=c\kappa$, for example when their counts always exceed $m/\kappa$.
The occurrence and impact of either condition in operational workloads remain
unmeasured.

The exact singleton counterexample nevertheless shows that the
published controlled-recovery claim for adjacent identifiers does not hold
for every allowed release. A revised claim or construction must address these
partial batches explicitly.
